\section{Implementation Details and Code Architecture}
\label{sec:implementation_details}

The MedRAGAgents codebase follows a modular Python architecture designed for clarity, extensibility, and maintainability. This section outlines key class implementations and interface abstractions.

\subsection{Base Agent Abstraction Layer (\texttt{agents/base\_agent.py})}

All agent modules inherit from \texttt{BaseAgent}, which wraps underlying API provider models (Google Gemini, OpenAI, Anthropic) with standard error handling, rate limiting, and output parsing.

Listing~\ref{lst:base_agent} demonstrates the base agent execution and regex extraction logic.

\begin{lstlisting}[language=Python, caption={BaseAgent Core Call and Answer Extraction Implementation.}, label={lst:base_agent}]
import re
import time
import google.generativeai as genai
import config as cfg

class BaseAgent:
    def __init__(self, llm_name: str = None, max_tokens: int = 1000):
        self.llm_name = llm_name or cfg.DEFAULT_LLM
        self.max_tokens = max_tokens
        self._init_client()

    def call(self, system: str, user: str) -> str:
        """Executes API invocation with retry logic."""
        for attempt in range(3):
            try:
                response = self.model.generate_content(
                    f"{system}\n\n{user}",
                    generation_config={"max_output_tokens": self.max_tokens}
                )
                return response.text.strip()
            except Exception as e:
                time.sleep(2 ** attempt)
        return ""

    @staticmethod
    def extract_answer_letter(text: str) -> str:
        """Extracts option letter (A-E) via regex fallback."""
        match = re.search(r'Option:\s*([A-E])', text, re.IGNORECASE)
        if match:
            return match.group(1).upper()
        match_alt = re.search(r'\b([A-E])\b', text)
        return match_alt.group(1).upper() if match_alt else ""
\end{lstlisting}

\subsection{Pipeline Orchestration (\texttt{pipeline.py})}

The pipeline orchestrator configures and executes system variants. Listing~\ref{lst:v3_pipeline} shows the core execution flow for Variant V3.

\begin{lstlisting}[language=Python, caption={Variant V3 Full Pipeline Execution Logic in pipeline.py.}, label={lst:v3_pipeline}]
class V3FullSystem:
    def run(self, question: str, options=None, gold_answer: str = "") -> dict:
        options_str = _format_options(options)
        rec = _empty_record(question, options, gold_answer)

        # Long-Term Memory Lookup
        cached = self.mem.long.get(question, options)
        if cached:
            rec.update(cached)
            rec["from_cache"] = True
            return rec

        # Short-Term Memory Reset
        self.mem.reset_short()

        # Step 1: Domain Classification
        q_domains = self.domain_agent.classify_question_domains(question)
        o_domains = self.domain_agent.classify_option_domains(question, options_str)
        self.mem.short.store("q_domains", q_domains)
        self.mem.short.store("o_domains", o_domains)

        # Step 2: Specialty Analyses
        q_analyses = self.analysis_agent.run_question_analyses(
            question, self.mem.short.get("q_domains")
        )
        o_analyses = self.analysis_agent.run_option_analyses(
            question, options_str, self.mem.short.get("o_domains"), q_analyses
        )

        # Step 3: MedRAG Retrieval & Context Injection
        rag_result = self.rag_agent.answer(question=question, options=options)
        if rag_result["snippets"]:
            q_analyses["RAG Evidence"] = self._summarise_snippets(rag_result["snippets"])

        # Step 4: Synthesis & Consensus Verification
        syn_report = self.synthesis_agent.synthesise(question, options_str, q_analyses, o_analyses)
        verif = self.verifier_agent.verify_and_answer(syn_report, q_domains + o_domains)

        rec.update({"pred_answer": verif["pred_answer"], "syn_report": verif["final_syn_report"]})
        self.mem.long.store(question, options, rec, autosave=True)
        return rec
\end{lstlisting}

\subsection{Memory Management Implementation (\texttt{memory.py})}

Listing~\ref{lst:long_term_memory} displays the SHA-256 key generation and caching implementation within \texttt{LongTermMemory}.

\begin{lstlisting}[language=Python, caption={SHA-256 Key Generation and Disk Caching in memory.py.}, label={lst:long_term_memory}]
import hashlib
import json
import os

class LongTermMemory:
    def __init__(self, cache_path: str = "./memory/long_term_cache.json"):
        self.cache_path = cache_path
        os.makedirs(os.path.dirname(cache_path), exist_ok=True)
        self._cache = self._load()

    @staticmethod
    def make_key(question: str, options: any) -> str:
        raw = question.strip() + str(options)
        return hashlib.sha256(raw.encode("utf-8")).hexdigest()[:16]

    def get(self, question: str, options: any):
        key = self.make_key(question, options)
        return self._cache.get(key)

    def store(self, question: str, options: any, record: dict, autosave: bool = True):
        key = self.make_key(question, options)
        self._cache[key] = record
        if autosave:
            self.save()
\end{lstlisting}
